\section*{Bloque A. Fundamentos teóricos y conceptuales}

\noindent\textit{Preguntas 1 a 30. Evalúan la comprensión rigurosa de los conceptos de teoría de juegos estáticos y dinámicos, estrategias dominantes, equilibrio de Nash, SPNE, modelos de oligopolio de Cournot, Bertrand y Stackelberg, estabilidad de cárteles y competencia monopolística.}

\q{En teoría de juegos, una estrategia de un jugador se define como estrictamente dominante si:}
\begin{opts}
    \item minimiza la pérdida máxima esperada que el jugador sufre en el peor escenario posible.
    \item proporciona un pago estrictamente mayor que cualquier otra estrategia ante toda acción rival.
    \item proporciona el pago máximo posible únicamente cuando todos los demás jugadores cooperan.
    \item iguala el pago esperado de todas las estrategias puras disponibles de los competidores.
\end{opts}

\q{Un perfil de estrategias en un juego estático con información completa constituye un Equilibrio de Nash si:}
\begin{opts}
    \item ningún jugador tiene incentivos para desviarse unilateralmente dadas las estrategias rivales.
    \item maximiza la suma agregada de los pagos conjuntos de todos los jugadores del mercado.
    \item todos los jugadores participantes obtienen exactamente el mismo nivel de pago esperado.
    \item cada jugador elige necesariamente una estrategia estrictamente dominante en estrategias puras.
\end{opts}

\q{En el juego clásico del Dilema del Prisionero en una sola etapa, el equilibrio de Nash:}
\begin{opts}
    \item se alcanza cuando ambos jugadores eligen no confesar o cooperar.
    \item no existe en estrategias puras, existiendo únicamente en estrategias mixtas.
    \item es eficiente en el sentido de Pareto, pues maximiza el bienestar social conjunto.
    \item es único, en estrategias dominantes, pero es ineficiente en el sentido de Pareto.
\end{opts}

\q{En un juego secuencial o dinámico representado en forma extensiva con información perfecta:}
\begin{opts}
    \item los jugadores desconocen las decisiones tomadas en etapas previas por los rivales.
    \item cada conjunto de información consta de un único nodo de decisión.
    \item los jugadores eligen sus acciones de forma simultánea e independiente.
    \item existen nodos de decisión que pertenecen a varios jugadores a la vez.
\end{opts}

\q{El método de inducción hacia atrás (backward induction) permite resolver juegos dinámicos finitos:}
\begin{opts}
    \item analizando primero los subjuegos en nodos terminales y retrocediendo hacia la raíz.
    \item maximizando la suma cooperativa de los pagos de todos los jugadores en la etapa inicial.
    \item eliminando de forma iterativa estrategias débilmente dominadas desde el nodo raíz.
    \item partiendo de las decisiones iniciales y proyectando pagos esperados hacia el futuro.
\end{opts}

\q{Un refinamiento fundamental del Equilibrio de Nash en juegos dinámicos es el Equilibrio de Nash Perfecto en Subjuegos (SPNE) propuesto por Reinhard Selten, el cual se caracteriza por:}
\begin{opts}
    \item admitir que los jugadores adopten amenazas inverosímiles para inducir cooperación tácita.
    \item descartar estrategias basadas en amenazas o promesas no creíbles en cualquier subjuego.
    \item aplicarse exclusivamente a juegos simultáneos finitos sin subjuegos propios relevantes.
    \item garantizar que el resultado coincida siempre con el equilibrio de colusión cooperativa.
\end{opts}

\q{En un juego estático finito donde un jugador adopta una estrategia mixta en equilibrio de Nash:}
\begin{opts}
    \item es indiferente entre todas las estrategias puras a las que asigna probabilidad positiva.
    \item obtiene un pago esperado estrictamente superior al de cualquiera de sus estrategias puras.
    \item asigna exactamente la misma probabilidad uniforme a todo el conjunto de acciones puras.
    \item asigna probabilidad igual a la unidad a la acción con menor varianza de pago esperado.
\end{opts}

\q{El Teorema Popular (Folk Theorem) para juegos repetidos infinitamente establece que:}
\begin{opts}
    \item cualquier acuerdo colusorio colapsa inevitablemente tras un número finito de periodos dados.
    \item si los jugadores son suficientemente pacientes, todo pago factible e individualmente racional es SPNE.
    \item la única solución posible en repetición infinita es repetir el equilibrio estático de una etapa.
    \item ningún resultado cooperativo puede sostenerse cuando los jugadores descuentan flujos futuros.
\end{opts}

\q{En el modelo clásico de oligopolio de Cournot, se asume que las empresas:}
\begin{opts}
    \item forman un cártel vinculante de precios y reparten el mercado mediante cuotas geográficas.
    \item eligen simultáneamente sus cantidades de producción tomando como dada la producción rival.
    \item compiten simultáneamente en precios homogéneos y la demanda residual ajusta las cantidades.
    \item eligen secuencialmente sus niveles de inversión en capacidad antes de fijar precios finales.
\end{opts}

\q{En el modelo de Cournot, la función de mejor respuesta (o función de reacción) de una empresa:}
\begin{opts}
    \item determina el precio óptimo de monopolio a cobrar ante cualquier tarifa anunciada por el rival.
    \item resulta independiente de los costos marginales propios si la curva de demanda es lineal.
    \item tiene pendiente positiva cuando las cantidades producidas actúan como sustitutos estratégicos.
    \item indica la cantidad que maximiza su beneficio para cada producción esperada de las rivales.
\end{opts}

\q{El modelo de oligopolio de Bertrand con bienes homogéneos y costos marginales constantes e idénticos asume que:}
\begin{opts}
    \item las empresas fijan precios escalonados según su capacidad instalada previa.
    \item los bienes presentan diferenciación horizontal espacial o de marca.
    \item la firma con menor precio captura la totalidad de la demanda del mercado.
    \item las empresas compiten fijando simultáneamente sus volúmenes de producción.
\end{opts}

\q{La denominada 'Paradoja de Bertrand' se refiere al resultado teórico según el cual:}
\begin{opts}
    \item un monopolista discriminador de precios genera exactamente el bienestar de competencia perfecta.
    \item la colusión perfecta resulta más estable conforme el número de competidores tiende a infinito.
    \item competir en precios con bien homogéneo lleva el precio al costo marginal y beneficio nulo.
    \item un cártel maximiza beneficios conjuntos cuando los costos unitarios de transporte son infinitos.
\end{opts}

\q{La solución propuesta por Francis Edgeworth para superar la paradoja de Bertrand se basa en incorporar:}
\begin{opts}
    \item asimetría de información bilateral sobre el costo fijo de entrada.
    \item restricciones rígidas de capacidad productiva en las empresas rivales.
    \item bienes perfectamente sustitutos producidos con costos marginales nulos.
    \item diferenciación vertical de calidad basada en inversión publicitaria.
\end{opts}

\q{Cuando dos empresas compiten en precios ofreciendo productos diferenciados (sustitutos imperfectos), sus funciones de mejor respuesta:}
\begin{opts}
    \item tienen pendiente negativa e idéntica magnitud que en el modelo de Cournot.
    \item son líneas perfectamente verticales e independientes de las decisiones rivales.
    \item presentan pendiente nula siempre que los costos marginales sean constantes.
    \item tienen pendiente positiva, reflejando complementariedad estratégica en precios.
\end{opts}

\q{Según la clasificación de Bulow, Geanakoplos y Klemperer (1985), las variables de decisión en oligopolio son 'sustitutos estratégicos' cuando:}
\begin{opts}
    \item un incremento en la acción de la empresa 1 reduce la respuesta óptima o acción de la empresa 2.
    \item un incremento en la acción de la empresa 1 induce a la empresa 2 a incrementar también su acción.
    \item las empresas fijan precios iguales a los costos medios mínimos de largo plazo.
    \item el producto carece de sustitutos cercanos en el mercado relevante.
\end{opts}

\q{En el modelo de liderazgo en cantidades de Stackelberg con una empresa líder y una seguidora:}
\begin{opts}
    \item el líder elige primero su volumen anticipando la función de reacción seguidora.
    \item ambas firmas obtienen exactamente los mismos beneficios que en Cournot simétrico.
    \item ambas empresas eligen su nivel de producción simultáneamente sin observarse.
    \item la seguidora actúa primero y el líder toma esa producción como un dato exógeno.
\end{opts}

\q{La ventaja del primer jugador (first-mover advantage) en el modelo de Stackelberg en cantidades requiere fundamentalmente:}
\begin{opts}
    \item que el mercado opere mediante una subasta simultánea a sobre cerrado de primer precio.
    \item que la empresa seguidora desconozca completamente el volumen elegido por la líder.
    \item que los costos marginales de la empresa líder superen estrictamente a los de su rival.
    \item que la empresa líder pueda comprometerse de manera creíble con su nivel de producción.
\end{opts}

\q{Considere una industria con demanda lineal y costos marginales constantes idénticos. El orden correcto de menor a mayor producción agregada ($Q$) es:}
\begin{opts}
    \item $Q_{Stackelberg} < Q_{Cournot} < Q_{Monopolio} < Q_{Competencia\ Perfecta}$
    \item $Q_{Monopolio} < Q_{Stackelberg} < Q_{Cournot} < Q_{Competencia\ Perfecta}$
    \item $Q_{Monopolio} < Q_{Cournot} < Q_{Stackelberg} < Q_{Competencia\ Perfecta}$
    \item $Q_{Cournot} < Q_{Monopolio} < Q_{Stackelberg} < Q_{Competencia\ Perfecta}$
\end{opts}

\q{Siguiendo el mismo mercado del caso anterior, el orden correcto de menor a mayor precio de mercado ($P$) es:}
\begin{opts}
    \item $P_{Competencia\ Perfecta} < P_{Monopolio} < P_{Cournot} < P_{Stackelberg}$
    \item $P_{Competencia\ Perfecta} < P_{Stackelberg} < P_{Cournot} < P_{Monopolio}$
    \item $P_{Monopolio} < P_{Cournot} < P_{Stackelberg} < P_{Competencia\ Perfecta}$
    \item $P_{Cournot} < P_{Stackelberg} < P_{Monopolio} < P_{Competencia\ Perfecta}$
\end{opts}

\q{La inestabilidad inherente de los acuerdos colusorios o cárteles en oligopolio se debe a que:}
\begin{opts}
    \item el beneficio conjunto del cártel resulta menor que la suma de beneficios de Cournot.
    \item el acuerdo colusorio incrementa sistemáticamente la elasticidad precio de la demanda.
    \item cada firma tiene incentivo unilateral a producir sobre su cuota si las demás cooperan.
    \item la normativa de defensa de la competencia incentiva la estabilidad de cárteles formales.
\end{opts}

\q{Bajo la estrategia del gatillo (grim trigger) en un juego repetido infinitamente, las empresas sostienen la colusión si:}
\begin{opts}
    \item la sanción por incumplimiento consiste en cobrar el precio monopólico por tres periodos.
    \item el número de competidores simétricos participantes dentro del acuerdo tiende a infinito.
    \item el valor presente de cooperar supera o iguala el desvío inmediato más el castigo perpetuo.
    \item el factor de descuento intertemporal $\delta$ cae por debajo de la tasa libre de riesgo.
\end{opts}

\q{En un oligopolio de Cournot simétrico con $n$ empresas que intentan coludir, a medida que el número de empresas $n$ aumenta:}
\begin{opts}
    \item el factor de descuento crítico $\delta^*$ sube, dificultando sostener el acuerdo colusorio.
    \item la ganancia inmediata de desviarse unilateralmente del acuerdo se reduce exactamente a cero.
    \item la cuota individual de mercado asignada dentro del cártel se expande proporcionalmente.
    \item el factor de descuento crítico $\delta^*$ disminuye, facilitando la estabilidad del cártel.
\end{opts}

\q{Uno de los supuestos esenciales del modelo de Competencia Monopolística formulado por Edward Chamberlin es:}
\begin{opts}
    \item la diferenciación de productos junto con libre entrada y salida a largo plazo.
    \item la interdependencia estratégica explícita entre un número reducido de oligopolistas.
    \item la presencia de elevadas barreras legales y tecnológicas a la entrada de rivales.
    \item la homogeneidad estricta de los bienes ofrecidos por todas las firmas competidoras.
\end{opts}

\q{El modelo de localización espacial de Hotelling formaliza el concepto de:}
\begin{opts}
    \item discriminación perfecta de precios de primer grado según disposición a pagar.
    \item diferenciación horizontal donde influye la preferencia o cercanía geográfica.
    \item monopolio natural multiproducto con funciones de costos subaditivas crecientes.
    \item diferenciación vertical donde todos coinciden en preferir mayor calidad objetiva.
\end{opts}

\q{A diferencia de la diferenciación horizontal, en un modelo de diferenciación vertical de productos:}
\begin{opts}
    \item las firmas compiten exclusivamente eligiendo cantidades simultáneas a la Cournot.
    \item a igual precio, todos los consumidores prefieren la variedad de mayor calidad.
    \item los costos fijos hundidos de entrada al mercado son necesariamente nulos.
    \item todos los consumidores poseen exactamente la misma disposición marginal a pagar.
\end{opts}

\q{En el equilibrio de largo plazo de una industria bajo competencia monopolística:}
\begin{opts}
    \item la demanda percibida es tangente al costo medio, resultando beneficio económico nulo.
    \item todas las firmas operan en el mínimo de su curva de costo total medio de largo plazo.
    \item el precio iguala exactamente al costo marginal y existen beneficios extraordinarios.
    \item el volumen agregado de mercado coincide con el equilibrio colusorio de monopolio.
\end{opts}

\q{El 'teorema del exceso de capacidad' en competencia monopolística establece que en el equilibrio de largo plazo:}
\begin{opts}
    \item las firmas producen por debajo de su escala mínima eficiente, con capacidad ociosa.
    \item el costo marginal de producción supera sistemáticamente al precio de equilibrio vigente.
    \item las firmas operan exactamente en la escala mínima eficiente minimizando su costo medio.
    \item las empresas producen por encima de su escala óptima, generando deseconomías de escala.
\end{opts}

\q{La evaluación del bienestar social en competencia monopolística frente a la competencia perfecta refleja una disyuntiva (trade-off) entre:}
\begin{opts}
    \item discriminación de precios de tercer grado y estabilidad macroeconómica sectorial.
    \item maximización de rentas monopólicas y recaudación por impuestos correctivos pigouvianos.
    \item reducción de costos salariales unitarios y eliminación de información asimétrica.
    \item eficiencia productiva de costo mínimo y valor social de mayor variedad de bienes.
\end{opts}

\q{En los mercados con competencia imperfecta, la publicidad informativa se diferencia de la publicidad persuasiva en que:}
\begin{opts}
    \item reduce costos de búsqueda difundiendo información objetiva sobre precios y atributos.
    \item disminuye la elasticidad precio de la demanda permitiendo márgenes más elevados.
    \item se encuentra prohibida por las agencias de competencia por inducir colusión tácita.
    \item incrementa las barreras estratégicas manipulando subjetivamente las preferencias.
\end{opts}

\q{En un duopolio de Cournot, la condición matemática que asegura la estabilidad dinámica del equilibrio de Nash es que:}
\begin{opts}
    \item el valor absoluto de la pendiente de cada función de reacción sea menor que la unidad.
    \item la suma de los costos marginales de todas las firmas participantes sea exactamente cero.
    \item las funciones de mejor respuesta se crucen necesariamente sobre el eje horizontal.
    \item el determinante de la matriz jacobiana de demandas inversas sea estrictamente negativo.
\end{opts}

